\documentclass[10pt,a4paper]{amsart}
\usepackage[margin=19mm]{geometry}
\usepackage{amsmath,amssymb,mathtools}
\usepackage[scr=rsfs,cal=euler]{mathalfa}
\usepackage{xfrac}
\usepackage[unicode,hidelinks,bookmarks=false]{hyperref}
\newcommand{\ie}{i.e.\ }
\newcommand{\cf}{cf.\ }
\newcommand{\resp}{respectively\ }
\newcommand{\Subsection}{\subsection}
\providecommand{\nobreakdash}{\nobreak\hbox{-}}
\setlength{\emergencystretch}{2em}
\multlinegap0pt
\usepackage{amssymb}
\usepackage{float}
\usepackage{tikz}
\usepackage{mathtools}
\newtheorem{theorem}{Theorem}
\title{On the Asymptotic Palindrome Density of Fibonacci Infinite Words}
\author{Duaa Abdullah, Jasem Hamoud}
\date{Source: arXiv:2511.06485v2 (CC BY 4.0)}
\begin{document}
\maketitle

\noindent\emph{Source and licence.} On the Asymptotic Palindrome Density of Fibonacci Infinite Words. Version arXiv:2511.06485v2 (CC BY 4.0), distributed by arXiv under the Creative Commons Attribution 4.0 International licence, http://creativecommons.org/licenses/by/4.0/, as recorded in the arXiv licence metadata for this exact version and shown on the abstract page https://arxiv.org/abs/2511.06485v2. Full licence notice: \texttt{licenses/arxiv-2511.06485-CC-BY-4.0.txt}.\par\smallskip

\noindent\textit{Editorial scope.} Counted unit: the article's theorem Thmmorsn1 (source0.tex lines 406--413). The article's complete original proof of it is delivered with it (source0.tex lines 415--440, one \texttt{proof} environment of 26 lines). No mathematics is written, repaired or re-derived: the statement and the proof are the article's own lines, in the article's own notation, and no retained line is substituted or re-flowed.  No further source-local context is needed: every cross-reference of every retained line resolves inside the two delivered spans.  Nothing was omitted from any retained span, so the package carries no omission record.  No retained line contains a citation, so the wrapper carries no bibliography and no bibliography entry.  No retained line contains a figure environment, an image-inclusion command or drawing code, so no image is delivered and the package declares no asset dependency.  Editorial formatting only, in addition: an AMS \texttt{amsart} preamble that restates the article's own declarations and macros used by the retained lines, namely the environments \texttt{theorem} on separate counters, together with the standard packages those lines need.  The retained environments print in this document as Theorem 1 in order of appearance, whereas the article prints the same items with the numbering of its own full document.  The complete native source of the article as distributed in the arXiv e-print for this exact version is bundled unchanged as \texttt{sources/188732/source0.tex}.
\begin{theorem}\label{Thmmorsn1}
The Thue--Morse sequence has balanced letter frequencies:
\begin{equation}\label{eq1Thmmorsn1}
\mathrm{dens}_0(t) = \mathrm{dens}_1(t) = \frac{1}{2},
\qquad\text{and hence}\qquad
\mathrm{dens}_0(t) + \mathrm{dens}_1(t) = 1.
\end{equation}
\end{theorem}

\begin{proof}
Let $n=2^k$ and consider the prefix $p_n$ of length $n$. We prove by induction on $k$ that
\[
|p_n|_0 = |p_n|_1 = 2^{k-1}.
\]
For $k=1$ we have $p_2=01$, so $|p_2|_0 = |p_2|_1 = 1$. Assume the claim holds for $k$, i.e., $|p_{2^k}|_0 = |p_{2^k}|_1 = 2^{k-1}$. The morphism description implies
\[
p_{2^{k+1}} = p_{2^k}\,\overline{p_{2^k}},
\]
where $\overline{w}$ denotes the bitwise complement of a binary word $w$. Thus
\[
|p_{2^{k+1}}|_0 = |p_{2^k}|_0 + |\overline{p_{2^k}}|_0 = |p_{2^k}|_0 + |p_{2^k}|_1 = 2^{k-1} + 2^{k-1} = 2^k,
\]
and similarly $|p_{2^{k+1}}|_1 = 2^k$. This proves the result for $n=2^k$.

For an arbitrary $n$, write $n=2^k + m$ with $0\le m<2^k$. Then
\[
p_n = p_{2^k}\, t_{2^k}\cdots t_{2^k+m-1},
\]
and the suffix $t_{2^k}\cdots t_{2^k+m-1}$ is the prefix of length $m$ of $\overline{t}$, since $t_{2^k+j}=1-t_j$. Using the exact balance at $2^k$ and the trivial bound $||p_m|_0-|p_m|_1|\le m$, one obtains
\[
\left|\frac{\lambda_n}{n} - \frac{1}{2}\right| \le \frac{m}{2^k+m},\qquad
\left|\frac{\alpha_n}{n} - \frac{1}{2}\right| \le \frac{m}{2^k+m},
\]
which tends to $0$ as $n\to\infty$. Hence the limits in~\eqref{eq1Thmmorsn1} exist and equal $1/2$.
\end{proof}

\end{document}
